% Ngân hàng câu hỏi Toán 11 (KNTT) -- Chương 2, Bài 5, Dạng 09
% Dạng: Dãy số: số hạng, cách cho dãy số, tính tăng giảm, bị chặn
% 100 câu: NB 30 -- TH 40 -- VD 20 -- VDC 10 ; mỗi mức: trắc nghiệm : đúng/sai : trả lời ngắn = 3 : 1 : 1
% Phương án ĐẦU TIÊN trong \cauTN luôn là đáp án đúng (gói trondethi tự đảo).
% Phần II: \y{nội dung}{d|s}. Phần III: đáp án đặt ở đối số cuối.

\cauTN[id=T11-C2-B5-D09-NB-001,muc=NB,loigiai={Thay $n=2$: $u_{2}=11$.}]{Cho dãy số $(u_n)$ với $u_n=3^n+2$. Số hạng $u_{2}$ bằng}
  {$11$.}
  {$29$.}{$5$.}{$8$.}

\cauTN[id=T11-C2-B5-D09-NB-002,muc=NB,loigiai={Tính lần lượt: $u_{1}=0, u_{2}=1, u_{3}=4$.}]{Cho dãy số $(u_n)$ xác định bởi $u_1=0$ và $u_{n+1}=3u_n+1$ với mọi $n\ge1$. Số hạng $u_{3}$ bằng}
  {$4$.}
  {$1$.}{$13$.}{$5$.}

\cauTN[id=T11-C2-B5-D09-NB-003,muc=NB,loigiai={Kiểm tra với $n=1,2,3,\ldots$ thấy $u_n=\dfrac{n}{n+1}$ phù hợp.}]{Cho dãy số có các số hạng đầu là $\dfrac12;\ \dfrac23;\ \dfrac34;\ \dfrac45;\ldots$ Số hạng tổng quát của dãy số này có thể là}
  {$u_n=\dfrac{n}{n+1}$.}
  {$u_n=\dfrac{n+1}{n}$.}{$u_n=\dfrac{n}{n+2}$.}{$u_n=\dfrac{1}{n+1}$.}

\cauTN[id=T11-C2-B5-D09-NB-004,muc=NB,loigiai={Thay $n=5$: $u_{5}=244$.}]{Cho dãy số $(u_n)$ với $u_n=3^n+1$. Số hạng $u_{5}$ bằng}
  {$244$.}
  {$730$.}{$82$.}{$16$.}

\cauTN[id=T11-C2-B5-D09-NB-005,muc=NB,loigiai={Tính lần lượt: $u_{1}=-1, u_{2}=-3, u_{3}=-7, u_{4}=-15$.}]{Cho dãy số $(u_n)$ xác định bởi $u_1=-1$ và $u_{n+1}=2u_n-1$ với mọi $n\ge1$. Số hạng $u_{4}$ bằng}
  {$-15$.}
  {$-7$.}{$-31$.}{$-16$.}

\cauTN[id=T11-C2-B5-D09-NB-006,muc=NB,loigiai={Kiểm tra với $n=1,2,3,\ldots$ thấy $u_n=2n-1$ phù hợp.}]{Cho dãy số có các số hạng đầu là $1;\ 3;\ 5;\ 7;\ 9;\ldots$ Số hạng tổng quát của dãy số này có thể là}
  {$u_n=2n-1$.}
  {$u_n=2n+1$.}{$u_n=n+2$.}{$u_n=n^2-1$.}

\cauTN[id=T11-C2-B5-D09-NB-007,muc=NB,loigiai={Thay $n=2$: $u_{2}=-3$.}]{Cho dãy số $(u_n)$ với $u_n=n^2-5n+3$. Số hạng $u_{2}$ bằng}
  {$-3$.}
  {$-4$.}{$-6$.}{$-1$.}

\cauTN[id=T11-C2-B5-D09-NB-008,muc=NB,loigiai={Tính lần lượt: $u_{1}=5, u_{2}=-14, u_{3}=24$.}]{Cho dãy số $(u_n)$ xác định bởi $u_1=5$ và $u_{n+1}=-2u_n-4$ với mọi $n\ge1$. Số hạng $u_{3}$ bằng}
  {$24$.}
  {$-14$.}{$-52$.}{$20$.}

\cauTN[id=T11-C2-B5-D09-NB-009,muc=NB,loigiai={Kiểm tra với $n=1,2,3,\ldots$ thấy $u_n=3n+2$ phù hợp.}]{Cho dãy số có các số hạng đầu là $5;\ 8;\ 11;\ 14;\ldots$ Số hạng tổng quát của dãy số này có thể là}
  {$u_n=3n+2$.}
  {$u_n=3n+5$.}{$u_n=5n$.}{$u_n=n+4$.}

\cauTN[id=T11-C2-B5-D09-NB-010,muc=NB,loigiai={Thay $n=7$: $u_{7}=13$.}]{Cho dãy số $(u_n)$ với $u_n=n^2-5n-1$. Số hạng $u_{7}$ bằng}
  {$13$.}
  {$5$.}{$23$.}{$14$.}

\cauTN[id=T11-C2-B5-D09-NB-011,muc=NB,loigiai={Tính lần lượt: $u_{1}=-2, u_{2}=-3, u_{3}=-\dfrac{7}{2}$.}]{Cho dãy số $(u_n)$ xác định bởi $u_1=-2$ và $u_{n+1}=\dfrac{1}{2}u_n-2$ với mọi $n\ge1$. Số hạng $u_{3}$ bằng}
  {$-\dfrac{7}{2}$.}
  {$-\dfrac{15}{4}$.}{$-\dfrac{11}{2}$.}{$-3$.}

\cauTN[id=T11-C2-B5-D09-NB-012,muc=NB,loigiai={Kiểm tra với $n=1,2,3,\ldots$ thấy $u_n=3^n$ phù hợp.}]{Cho dãy số có các số hạng đầu là $3;\ 9;\ 27;\ 81;\ldots$ Số hạng tổng quát của dãy số này có thể là}
  {$u_n=3^n$.}
  {$u_n=3n$.}{$u_n=n^3$.}{$u_n=3^{n+1}$.}

\cauTN[id=T11-C2-B5-D09-NB-013,muc=NB,loigiai={Thay $n=3$: $u_{3}=-\dfrac{5}{2}$.}]{Cho dãy số $(u_n)$ với $u_n=\dfrac{(-1)^n(3n+1)}{n+1}$. Số hạng $u_{3}$ bằng}
  {$-\dfrac{5}{2}$.}
  {$\dfrac{5}{2}$.}{$-\dfrac{3}{2}$.}{$-\dfrac{10}{3}$.}

\cauTN[id=T11-C2-B5-D09-NB-014,muc=NB,loigiai={Tính lần lượt: $u_{1}=0, u_{2}=0, u_{3}=0$.}]{Cho dãy số $(u_n)$ xác định bởi $u_1=0$ và $u_{n+1}=2u_n$ với mọi $n\ge1$. Số hạng $u_{3}$ bằng}
  {$0$.}
  {$-1$.}{$1$.}{$2$.}

\cauTN[id=T11-C2-B5-D09-NB-015,muc=NB,loigiai={Kiểm tra với $n=1,2,3,\ldots$ thấy $u_n=n^2$ phù hợp.}]{Cho dãy số có các số hạng đầu là $1;\ 4;\ 9;\ 16;\ 25;\ldots$ Số hạng tổng quát của dãy số này có thể là}
  {$u_n=n^2$.}
  {$u_n=3n-2$.}{$u_n=2n-1$.}{$u_n=n^2+1$.}

\cauTN[id=T11-C2-B5-D09-NB-016,muc=NB,loigiai={Thay $n=2$: $u_{2}=-\dfrac{2}{7}$.}]{Cho dãy số $(u_n)$ với $u_n=\dfrac{-n}{2n+3}$. Số hạng $u_{2}$ bằng}
  {$-\dfrac{2}{7}$.}
  {$\dfrac{2}{7}$.}{$-\dfrac{1}{5}$.}{$-\dfrac{1}{3}$.}

\cauTN[id=T11-C2-B5-D09-NB-017,muc=NB,loigiai={Tính lần lượt: $u_{1}=-2, u_{2}=-7, u_{3}=-17, u_{4}=-37, u_{5}=-77$.}]{Cho dãy số $(u_n)$ xác định bởi $u_1=-2$ và $u_{n+1}=2u_n-3$ với mọi $n\ge1$. Số hạng $u_{5}$ bằng}
  {$-77$.}
  {$-37$.}{$-157$.}{$-80$.}

\cauTN[id=T11-C2-B5-D09-NB-018,muc=NB,loigiai={Kiểm tra với $n=1,2,3,\ldots$ thấy $u_n=2^n$ phù hợp.}]{Cho dãy số có các số hạng đầu là $2;\ 4;\ 8;\ 16;\ 32;\ldots$ Số hạng tổng quát của dãy số này có thể là}
  {$u_n=2^n$.}
  {$u_n=2n$.}{$u_n=n^2$.}{$u_n=2^{n-1}$.}

\cauDS[id=T11-C2-B5-D09-NB-019,muc=NB,loigiai={$u_{n+1}-u_n=\dfrac{7}{(n+1)(n+2)}$ > $0$; $u_n=4-\dfrac{7}{n+1}$. a) Đúng. b) Đúng. c) Đúng. d) Đúng.}]{Cho dãy số $(u_n)$ với $u_n=\dfrac{4n-3}{n+1}$.}{%
  \y{$u_1=\dfrac{1}{2}$.}{d}
  \y{$u_{2}=\dfrac{5}{3}$.}{d}
  \y{$(u_n)$ là dãy số tăng.}{d}
  \y{$u_n<4$ với mọi $n\in\mathbb{N}^*$.}{d}}

\cauDS[id=T11-C2-B5-D09-NB-020,muc=NB,loigiai={$u_{n+1}-u_n=\dfrac{-2}{(n+1)(n+2)}$ < $0$; $u_n=4+\dfrac{2}{n+1}$. a) Đúng. b) Đúng. c) Sai. d) Sai.}]{Cho dãy số $(u_n)$ với $u_n=\dfrac{4n+6}{n+1}$.}{%
  \y{$u_1=5$.}{d}
  \y{$u_{5}=\dfrac{13}{3}$.}{d}
  \y{$(u_n)$ là dãy số tăng.}{s}
  \y{$u_n<4$ với mọi $n\in\mathbb{N}^*$.}{s}}

\cauDS[id=T11-C2-B5-D09-NB-021,muc=NB,loigiai={$u_{n+1}-u_n=\dfrac{-13}{(n+2)(n+3)}$ < $0$; $u_n=-5+\dfrac{13}{n+2}$. a) Đúng. b) Sai. c) Sai. d) Sai.}]{Cho dãy số $(u_n)$ với $u_n=\dfrac{-5n+3}{n+2}$.}{%
  \y{$u_1=-\dfrac{2}{3}$.}{d}
  \y{$u_{3}=-\dfrac{7}{5}$.}{s}
  \y{$(u_n)$ là dãy số tăng.}{s}
  \y{$u_n<-5$ với mọi $n\in\mathbb{N}^*$.}{s}}

\cauDS[id=T11-C2-B5-D09-NB-022,muc=NB,loigiai={$u_{n+1}-u_n=\dfrac{-23}{(n+5)(n+6)}$ < $0$; $u_n=-5+\dfrac{23}{n+5}$. a) Đúng. b) Đúng. c) Sai. d) Sai.}]{Cho dãy số $(u_n)$ với $u_n=\dfrac{-5n-2}{n+5}$.}{%
  \y{$u_1=-\dfrac{7}{6}$.}{d}
  \y{$u_{5}=-\dfrac{27}{10}$.}{d}
  \y{$(u_n)$ là dãy số tăng.}{s}
  \y{$u_n<-5$ với mọi $n\in\mathbb{N}^*$.}{s}}

\cauDS[id=T11-C2-B5-D09-NB-023,muc=NB,loigiai={$u_{n+1}-u_n=\dfrac{21}{(n+4)(n+5)}$ > $0$; $u_n=4-\dfrac{21}{n+4}$. a) Đúng. b) Sai. c) Đúng. d) Đúng.}]{Cho dãy số $(u_n)$ với $u_n=\dfrac{4n-5}{n+4}$.}{%
  \y{$u_1=-\dfrac{1}{5}$.}{d}
  \y{$u_{2}=\dfrac{3}{2}$.}{s}
  \y{$(u_n)$ là dãy số tăng.}{d}
  \y{$u_n<4$ với mọi $n\in\mathbb{N}^*$.}{d}}

\cauDS[id=T11-C2-B5-D09-NB-024,muc=NB,loigiai={$u_{n+1}-u_n=\dfrac{-2}{(n+1)(n+2)}$ < $0$; $u_n=-4+\dfrac{2}{n+1}$. a) Đúng. b) Đúng. c) Sai. d) Sai.}]{Cho dãy số $(u_n)$ với $u_n=\dfrac{-4n-2}{n+1}$.}{%
  \y{$u_1=-3$.}{d}
  \y{$u_{4}=-\dfrac{18}{5}$.}{d}
  \y{$(u_n)$ là dãy số tăng.}{s}
  \y{$u_n<-4$ với mọi $n\in\mathbb{N}^*$.}{s}}

\cauTLN[id=T11-C2-B5-D09-NB-025,muc=NB,loigiai={$u_{3}=23$.}]{Cho dãy số $(u_n)$ với $u_n=3n^2-4$. Tính $u_{3}$.}{23}

\cauTLN[id=T11-C2-B5-D09-NB-026,muc=NB,loigiai={$u_{6}=37$.}]{Cho dãy số $(u_n)$ với $u_n=2n^2-5n-5$. Tính $u_{6}$.}{37}

\cauTLN[id=T11-C2-B5-D09-NB-027,muc=NB,loigiai={Giải $n^2-5n+2=16$ với $n\in\mathbb{N}^*$ được $n=7$.}]{Cho dãy số $(u_n)$ với $u_n=n^2-5n+2$. Số $16$ là số hạng thứ mấy của dãy?}{7}

\cauTLN[id=T11-C2-B5-D09-NB-028,muc=NB,loigiai={$u_{8}=153$.}]{Cho dãy số $(u_n)$ với $u_n=3n^2-4n-7$. Tính $u_{8}$.}{153}

\cauTLN[id=T11-C2-B5-D09-NB-029,muc=NB,loigiai={Giải $n^2+3n=208$ với $n\in\mathbb{N}^*$ được $n=13$.}]{Cho dãy số $(u_n)$ với $u_n=n^2+3n$. Số $208$ là số hạng thứ mấy của dãy?}{13}

\cauTLN[id=T11-C2-B5-D09-NB-030,muc=NB,loigiai={$u_{7}=142$.}]{Cho dãy số $(u_n)$ với $u_n=3n^2-2n+9$. Tính $u_{7}$.}{142}

\cauTN[id=T11-C2-B5-D09-TH-031,muc=TH,loigiai={Xét hiệu $u_{n+1}-u_n$ (hoặc tỉ số $\dfrac{u_{n+1}}{u_n}$ với dãy dương): dãy $u_n=\dfrac{2n+1}{n+3}$ là dãy tăng.}]{Trong các dãy số $(u_n)$ sau, dãy số nào là dãy số tăng?}
  {$u_n=\dfrac{2n+1}{n+3}$.}
  {$u_n=\dfrac{3}{n+2}$.}{$u_n=(-1)^n\cdot n$.}{$u_n=n^2-4n$.}

\cauTN[id=T11-C2-B5-D09-TH-032,muc=TH,loigiai={$u_n=-3+\dfrac{8}{n+4}$ nên dãy giảm và $-3<u_n\le-\dfrac{7}{5}$.}]{Cho dãy số $(u_n)$ với $u_n=\dfrac{-3n-4}{n+4}$. Khẳng định nào sau đây đúng?}
  {$(u_n)$ là dãy số giảm và bị chặn.}
  {$(u_n)$ là dãy số tăng và bị chặn.}{$(u_n)$ là dãy số giảm và không bị chặn trên.}{$(u_n)$ không tăng, không giảm.}

\cauTN[id=T11-C2-B5-D09-TH-033,muc=TH,loigiai={Tính lần lượt: $u_{1}=-2, u_{2}=-9, u_{3}=-30, u_{4}=-93$.}]{Cho dãy số $(u_n)$ xác định bởi $u_1=-2$ và $u_{n+1}=3u_n-3$ với mọi $n\ge1$. Số hạng $u_{4}$ bằng}
  {$-93$.}
  {$-96$.}{$-30$.}{$-282$.}

\cauTN[id=T11-C2-B5-D09-TH-034,muc=TH,loigiai={Xét hiệu $u_{n+1}-u_n$ (hoặc tỉ số $\dfrac{u_{n+1}}{u_n}$ với dãy dương): dãy $u_n=\sqrt{n+1}$ là dãy tăng.}]{Trong các dãy số $(u_n)$ sau, dãy số nào là dãy số tăng?}
  {$u_n=\sqrt{n+1}$.}
  {$u_n=\dfrac{3}{n+2}$.}{$u_n=(-1)^n\cdot n$.}{$u_n=\dfrac{2-n}{n+1}$.}

\cauTN[id=T11-C2-B5-D09-TH-035,muc=TH,loigiai={Tính lần lượt: $u_{1}=-3, u_{2}=-\dfrac{9}{2}, u_{3}=-\dfrac{21}{4}, u_{4}=-\dfrac{45}{8}, u_{5}=-\dfrac{93}{16}$.}]{Cho dãy số $(u_n)$ xác định bởi $u_1=-3$ và $u_{n+1}=\dfrac{1}{2}u_n-3$ với mọi $n\ge1$. Số hạng $u_{5}$ bằng}
  {$-\dfrac{93}{16}$.}
  {$-\dfrac{189}{32}$.}{$-\dfrac{45}{8}$.}{$-\dfrac{141}{16}$.}

\cauTN[id=T11-C2-B5-D09-TH-036,muc=TH,loigiai={Xét hiệu $u_{n+1}-u_n$ (hoặc tỉ số $\dfrac{u_{n+1}}{u_n}$ với dãy dương): dãy $u_n=\dfrac{n+5}{n+1}$ là dãy giảm.}]{Trong các dãy số $(u_n)$ sau, dãy số nào là dãy số giảm?}
  {$u_n=\dfrac{n+5}{n+1}$.}
  {$u_n=n^2-4n$.}{$u_n=\dfrac{2n+1}{n+3}$.}{$u_n=(-1)^n\cdot n$.}

\cauTN[id=T11-C2-B5-D09-TH-037,muc=TH,loigiai={Tính lần lượt: $u_{1}=5, u_{2}=13, u_{3}=29$.}]{Cho dãy số $(u_n)$ xác định bởi $u_1=5$ và $u_{n+1}=2u_n+3$ với mọi $n\ge1$. Số hạng $u_{3}$ bằng}
  {$29$.}
  {$13$.}{$61$.}{$32$.}

\cauTN[id=T11-C2-B5-D09-TH-038,muc=TH,loigiai={Xét hiệu $u_{n+1}-u_n$ (hoặc tỉ số $\dfrac{u_{n+1}}{u_n}$ với dãy dương): dãy $u_n=\dfrac{n}{n^2+1}$ là dãy giảm.}]{Trong các dãy số $(u_n)$ sau, dãy số nào là dãy số giảm?}
  {$u_n=\dfrac{n}{n^2+1}$.}
  {$u_n=3^n-1$.}{$u_n=\dfrac{4n-1}{n+2}$.}{$u_n=\dfrac{2n+1}{n+3}$.}

\cauTN[id=T11-C2-B5-D09-TH-039,muc=TH,loigiai={$u_n=-1+\dfrac{1}{n+2}$ nên dãy giảm và $-1<u_n\le-\dfrac{2}{3}$.}]{Cho dãy số $(u_n)$ với $u_n=\dfrac{-n-1}{n+2}$. Khẳng định nào sau đây đúng?}
  {$(u_n)$ là dãy số giảm và bị chặn.}
  {$(u_n)$ là dãy số tăng và bị chặn.}{$(u_n)$ là dãy số giảm và không bị chặn trên.}{$(u_n)$ không tăng, không giảm.}

\cauTN[id=T11-C2-B5-D09-TH-040,muc=TH,loigiai={Tính lần lượt: $u_{1}=0, u_{2}=2, u_{3}=-2, u_{4}=6, u_{5}=-10$.}]{Cho dãy số $(u_n)$ xác định bởi $u_1=0$ và $u_{n+1}=-2u_n+2$ với mọi $n\ge1$. Số hạng $u_{5}$ bằng}
  {$-10$.}
  {$22$.}{$-8$.}{$6$.}

\cauTN[id=T11-C2-B5-D09-TH-041,muc=TH,loigiai={Xét hiệu $u_{n+1}-u_n$ (hoặc tỉ số $\dfrac{u_{n+1}}{u_n}$ với dãy dương): dãy $u_n=\dfrac{2-n}{n+1}$ là dãy giảm.}]{Trong các dãy số $(u_n)$ sau, dãy số nào là dãy số giảm?}
  {$u_n=\dfrac{2-n}{n+1}$.}
  {$u_n=\dfrac{2n+1}{n+3}$.}{$u_n=\sin\dfrac{n\pi}{2}$.}{$u_n=(-1)^n\cdot n$.}

\cauTN[id=T11-C2-B5-D09-TH-042,muc=TH,loigiai={Tính lần lượt: $u_{1}=5, u_{2}=\dfrac{5}{2}, u_{3}=\dfrac{5}{4}, u_{4}=\dfrac{5}{8}$.}]{Cho dãy số $(u_n)$ xác định bởi $u_1=5$ và $u_{n+1}=\dfrac{1}{2}u_n$ với mọi $n\ge1$. Số hạng $u_{4}$ bằng}
  {$\dfrac{5}{8}$.}
  {$\dfrac{13}{8}$.}{$\dfrac{5}{16}$.}{$-\dfrac{3}{8}$.}

\cauTN[id=T11-C2-B5-D09-TH-043,muc=TH,loigiai={Xét hiệu $u_{n+1}-u_n$ (hoặc tỉ số $\dfrac{u_{n+1}}{u_n}$ với dãy dương): dãy $u_n=\dfrac{4n-1}{n+2}$ là dãy tăng.}]{Trong các dãy số $(u_n)$ sau, dãy số nào là dãy số tăng?}
  {$u_n=\dfrac{4n-1}{n+2}$.}
  {$u_n=5-2n$.}{$u_n=\dfrac{n+5}{n+1}$.}{$u_n=\sin\dfrac{n\pi}{2}$.}

\cauTN[id=T11-C2-B5-D09-TH-044,muc=TH,loigiai={Tính lần lượt: $u_{1}=-3, u_{2}=-4, u_{3}=-6, u_{4}=-10$.}]{Cho dãy số $(u_n)$ xác định bởi $u_1=-3$ và $u_{n+1}=2u_n+2$ với mọi $n\ge1$. Số hạng $u_{4}$ bằng}
  {$-10$.}
  {$-18$.}{$-8$.}{$-6$.}

\cauTN[id=T11-C2-B5-D09-TH-045,muc=TH,loigiai={Tính lần lượt: $u_{1}=5, u_{2}=19, u_{3}=61$.}]{Cho dãy số $(u_n)$ xác định bởi $u_1=5$ và $u_{n+1}=3u_n+4$ với mọi $n\ge1$. Số hạng $u_{3}$ bằng}
  {$61$.}
  {$65$.}{$187$.}{$19$.}

\cauTN[id=T11-C2-B5-D09-TH-046,muc=TH,loigiai={Tính lần lượt: $u_{1}=2, u_{2}=7, u_{3}=22, u_{4}=67, u_{5}=202$.}]{Cho dãy số $(u_n)$ xác định bởi $u_1=2$ và $u_{n+1}=3u_n+1$ với mọi $n\ge1$. Số hạng $u_{5}$ bằng}
  {$202$.}
  {$67$.}{$203$.}{$607$.}

\cauTN[id=T11-C2-B5-D09-TH-047,muc=TH,loigiai={$u_n=4-\dfrac{22}{n+5}$ nên dãy tăng và $\dfrac{1}{3}\le u_n<4$.}]{Cho dãy số $(u_n)$ với $u_n=\dfrac{4n-2}{n+5}$. Khẳng định nào sau đây đúng?}
  {$(u_n)$ là dãy số tăng và bị chặn.}
  {$(u_n)$ là dãy số giảm và bị chặn.}{$(u_n)$ là dãy số tăng và không bị chặn trên.}{$(u_n)$ không tăng, không giảm.}

\cauTN[id=T11-C2-B5-D09-TH-048,muc=TH,loigiai={Tính lần lượt: $u_{1}=5, u_{2}=-7, u_{3}=17$.}]{Cho dãy số $(u_n)$ xác định bởi $u_1=5$ và $u_{n+1}=-2u_n+3$ với mọi $n\ge1$. Số hạng $u_{3}$ bằng}
  {$17$.}
  {$20$.}{$-7$.}{$-31$.}

\cauTN[id=T11-C2-B5-D09-TH-049,muc=TH,loigiai={Tính lần lượt: $u_{1}=-1, u_{2}=0, u_{3}=3, u_{4}=12$.}]{Cho dãy số $(u_n)$ xác định bởi $u_1=-1$ và $u_{n+1}=3u_n+3$ với mọi $n\ge1$. Số hạng $u_{4}$ bằng}
  {$12$.}
  {$3$.}{$39$.}{$15$.}

\cauTN[id=T11-C2-B5-D09-TH-050,muc=TH,loigiai={Xét hiệu $u_{n+1}-u_n$ (hoặc tỉ số $\dfrac{u_{n+1}}{u_n}$ với dãy dương): dãy $u_n=\dfrac{3}{n+2}$ là dãy giảm.}]{Trong các dãy số $(u_n)$ sau, dãy số nào là dãy số giảm?}
  {$u_n=\dfrac{3}{n+2}$.}
  {$u_n=\dfrac{2n+1}{n+3}$.}{$u_n=\left(-\dfrac12\right)^n$.}{$u_n=3^n-1$.}

\cauTN[id=T11-C2-B5-D09-TH-051,muc=TH,loigiai={Tính lần lượt: $u_{1}=2, u_{2}=0, u_{3}=-1, u_{4}=-\dfrac{3}{2}, u_{5}=-\dfrac{7}{4}$.}]{Cho dãy số $(u_n)$ xác định bởi $u_1=2$ và $u_{n+1}=\dfrac{1}{2}u_n-1$ với mọi $n\ge1$. Số hạng $u_{5}$ bằng}
  {$-\dfrac{7}{4}$.}
  {$-\dfrac{3}{2}$.}{$-\dfrac{11}{4}$.}{$-\dfrac{15}{8}$.}

\cauTN[id=T11-C2-B5-D09-TH-052,muc=TH,loigiai={Tính lần lượt: $u_{1}=-2, u_{2}=-3, u_{3}=-5$.}]{Cho dãy số $(u_n)$ xác định bởi $u_1=-2$ và $u_{n+1}=2u_n+1$ với mọi $n\ge1$. Số hạng $u_{3}$ bằng}
  {$-5$.}
  {$-9$.}{$-3$.}{$-4$.}

\cauTN[id=T11-C2-B5-D09-TH-053,muc=TH,loigiai={Tính lần lượt: $u_{1}=3, u_{2}=-2, u_{3}=8, u_{4}=-12, u_{5}=28$.}]{Cho dãy số $(u_n)$ xác định bởi $u_1=3$ và $u_{n+1}=-2u_n+4$ với mọi $n\ge1$. Số hạng $u_{5}$ bằng}
  {$28$.}
  {$-52$.}{$32$.}{$-12$.}

\cauTN[id=T11-C2-B5-D09-TH-054,muc=TH,loigiai={Xét hiệu $u_{n+1}-u_n$ (hoặc tỉ số $\dfrac{u_{n+1}}{u_n}$ với dãy dương): dãy $u_n=3^n-1$ là dãy tăng.}]{Trong các dãy số $(u_n)$ sau, dãy số nào là dãy số tăng?}
  {$u_n=3^n-1$.}
  {$u_n=5-2n$.}{$u_n=\dfrac{n}{n^2+1}$.}{$u_n=\dfrac{2-n}{n+1}$.}

\cauDS[id=T11-C2-B5-D09-TH-055,muc=TH,loigiai={$u_{n+1}-u_n=\dfrac{-10}{(n+4)(n+5)}$ < $0$; $u_n=-3+\dfrac{10}{n+4}$. a) Đúng. b) Đúng. c) Sai. d) Sai.}]{Cho dãy số $(u_n)$ với $u_n=\dfrac{-3n-2}{n+4}$.}{%
  \y{$u_1=-1$.}{d}
  \y{$u_{2}=-\dfrac{4}{3}$.}{d}
  \y{$(u_n)$ là dãy số tăng.}{s}
  \y{$u_n<-3$ với mọi $n\in\mathbb{N}^*$.}{s}}

\cauDS[id=T11-C2-B5-D09-TH-056,muc=TH,loigiai={$u_{n+1}-u_n=\dfrac{10}{(n+3)(n+4)}$ > $0$; $u_n=5-\dfrac{10}{n+3}$. a) Đúng. b) Đúng. c) Đúng. d) Đúng.}]{Cho dãy số $(u_n)$ với $u_n=\dfrac{5n+5}{n+3}$.}{%
  \y{$u_1=\dfrac{5}{2}$.}{d}
  \y{$u_{3}=\dfrac{10}{3}$.}{d}
  \y{$(u_n)$ là dãy số tăng.}{d}
  \y{$u_n<5$ với mọi $n\in\mathbb{N}^*$.}{d}}

\cauDS[id=T11-C2-B5-D09-TH-057,muc=TH,loigiai={$u_{n+1}-u_n=\dfrac{4}{(n+2)(n+3)}$ > $0$; $u_n=-1-\dfrac{4}{n+2}$. a) Đúng. b) Sai. c) Đúng. d) Đúng.}]{Cho dãy số $(u_n)$ với $u_n=\dfrac{-n-6}{n+2}$.}{%
  \y{$u_1=-\dfrac{7}{3}$.}{d}
  \y{$u_{2}=-1$.}{s}
  \y{$(u_n)$ là dãy số tăng.}{d}
  \y{$u_n<-1$ với mọi $n\in\mathbb{N}^*$.}{d}}

\cauDS[id=T11-C2-B5-D09-TH-058,muc=TH,loigiai={$u_{n+1}-u_n=\dfrac{9}{(n+1)(n+2)}$ > $0$; $u_n=5-\dfrac{9}{n+1}$. a) Đúng. b) Đúng. c) Đúng. d) Đúng.}]{Cho dãy số $(u_n)$ với $u_n=\dfrac{5n-4}{n+1}$.}{%
  \y{$u_1=\dfrac{1}{2}$.}{d}
  \y{$u_{4}=\dfrac{16}{5}$.}{d}
  \y{$(u_n)$ là dãy số tăng.}{d}
  \y{$u_n<5$ với mọi $n\in\mathbb{N}^*$.}{d}}

\cauDS[id=T11-C2-B5-D09-TH-059,muc=TH,loigiai={$u_{n+1}-u_n=\dfrac{17}{(n+4)(n+5)}$ > $0$; $u_n=3-\dfrac{17}{n+4}$. a) Đúng. b) Sai. c) Đúng. d) Đúng.}]{Cho dãy số $(u_n)$ với $u_n=\dfrac{3n-5}{n+4}$.}{%
  \y{$u_1=-\dfrac{2}{5}$.}{d}
  \y{$u_{5}=\dfrac{19}{9}$.}{s}
  \y{$(u_n)$ là dãy số tăng.}{d}
  \y{$u_n<3$ với mọi $n\in\mathbb{N}^*$.}{d}}

\cauDS[id=T11-C2-B5-D09-TH-060,muc=TH,loigiai={$u_{n+1}-u_n=\dfrac{-3}{(n+2)(n+3)}$ < $0$; $u_n=-4+\dfrac{3}{n+2}$. a) Đúng. b) Đúng. c) Sai. d) Sai.}]{Cho dãy số $(u_n)$ với $u_n=\dfrac{-4n-5}{n+2}$.}{%
  \y{$u_1=-3$.}{d}
  \y{$u_{2}=-\dfrac{13}{4}$.}{d}
  \y{$(u_n)$ là dãy số tăng.}{s}
  \y{$u_n<-4$ với mọi $n\in\mathbb{N}^*$.}{s}}

\cauDS[id=T11-C2-B5-D09-TH-061,muc=TH,loigiai={$u_{n+1}-u_n=\dfrac{-4}{(n+1)(n+2)}$ < $0$; $u_n=-1+\dfrac{4}{n+1}$. a) Đúng. b) Sai. c) Sai. d) Sai.}]{Cho dãy số $(u_n)$ với $u_n=\dfrac{-n+3}{n+1}$.}{%
  \y{$u_1=1$.}{d}
  \y{$u_{4}=\dfrac{4}{5}$.}{s}
  \y{$(u_n)$ là dãy số tăng.}{s}
  \y{$u_n<-1$ với mọi $n\in\mathbb{N}^*$.}{s}}

\cauDS[id=T11-C2-B5-D09-TH-062,muc=TH,loigiai={$u_{n+1}-u_n=\dfrac{-16}{(n+3)(n+4)}$ < $0$; $u_n=-4+\dfrac{16}{n+3}$. a) Đúng. b) Đúng. c) Sai. d) Sai.}]{Cho dãy số $(u_n)$ với $u_n=\dfrac{-4n+4}{n+3}$.}{%
  \y{$u_1=0$.}{d}
  \y{$u_{2}=-\dfrac{4}{5}$.}{d}
  \y{$(u_n)$ là dãy số tăng.}{s}
  \y{$u_n<-4$ với mọi $n\in\mathbb{N}^*$.}{s}}

\cauTLN[id=T11-C2-B5-D09-TH-063,muc=TH,loigiai={Giải $n^2-n+8=64$ với $n\in\mathbb{N}^*$ được $n=8$.}]{Cho dãy số $(u_n)$ với $u_n=n^2-n+8$. Số $64$ là số hạng thứ mấy của dãy?}{8}

\cauTLN[id=T11-C2-B5-D09-TH-064,muc=TH,loigiai={$u_{3}=8$.}]{Cho dãy số $(u_n)$ với $u_n=2n^2-6n+8$. Tính $u_{3}$.}{8}

\cauTLN[id=T11-C2-B5-D09-TH-065,muc=TH,loigiai={Giải $n^2+4n+8=85$ với $n\in\mathbb{N}^*$ được $n=7$.}]{Cho dãy số $(u_n)$ với $u_n=n^2+4n+8$. Số $85$ là số hạng thứ mấy của dãy?}{7}

\cauTLN[id=T11-C2-B5-D09-TH-066,muc=TH,loigiai={$u_{5}=40$.}]{Cho dãy số $(u_n)$ với $u_n=3n^2-6n-5$. Tính $u_{5}$.}{40}

\cauTLN[id=T11-C2-B5-D09-TH-067,muc=TH,loigiai={Giải $n^2-4n+3=120$ với $n\in\mathbb{N}^*$ được $n=13$.}]{Cho dãy số $(u_n)$ với $u_n=n^2-4n+3$. Số $120$ là số hạng thứ mấy của dãy?}{13}

\cauTLN[id=T11-C2-B5-D09-TH-068,muc=TH,loigiai={$u_{4}=36$.}]{Cho dãy số $(u_n)$ với $u_n=3n^2-2n-4$. Tính $u_{4}$.}{36}

\cauTLN[id=T11-C2-B5-D09-TH-069,muc=TH,loigiai={Giải $n^2-2n-2=78$ với $n\in\mathbb{N}^*$ được $n=10$.}]{Cho dãy số $(u_n)$ với $u_n=n^2-2n-2$. Số $78$ là số hạng thứ mấy của dãy?}{10}

\cauTLN[id=T11-C2-B5-D09-TH-070,muc=TH,loigiai={$u_{6}=58$.}]{Cho dãy số $(u_n)$ với $u_n=n^2+3n+4$. Tính $u_{6}$.}{58}

\cauTN[id=T11-C2-B5-D09-VD-071,muc=VD,loigiai={$u_1=4$, $u_{n+1}=u_n+3$ nên $u_n=3n+1$; $u_{6}=19$.}]{Bạn Nam dùng các que diêm giống nhau xếp thành một dãy các hình vuông liền nhau trên mặt bàn: hình thứ nhất là $1$ hình vuông cần $4$ que; mỗi lần xếp thêm một hình vuông kề bên phải hình vừa xếp, hai hình vuông dùng chung một cạnh nên chỉ cần thêm $3$ que. Gọi $u_n$ là số que diêm cần dùng để xếp được dãy gồm $n$ hình vuông. Số que diêm cần dùng để xếp dãy gồm $6$ hình vuông là}
  {$19$.}
  {$20$.}{$17$.}{$24$.}

\cauTN[id=T11-C2-B5-D09-VD-072,muc=VD,loigiai={$u_n=5+2n$; $u_{12}=29$.}]{Đầu năm học, bạn Việt có sẵn $5$ trăm nghìn đồng trong ống heo tiết kiệm. Kể từ tháng thứ nhất, cuối mỗi tháng bạn bỏ thêm vào ống heo $2$ trăm nghìn đồng tiền tiết kiệm từ tiền ăn sáng và không rút ra lần nào. Gọi $u_n$ (trăm nghìn đồng) là số tiền có trong ống heo sau $n$ tháng. Sau $12$ tháng, số tiền trong ống heo là}
  {$29$ trăm nghìn đồng.}
  {$27$ trăm nghìn đồng.}{$24$ trăm nghìn đồng.}{$31$ trăm nghìn đồng.}

\cauTN[id=T11-C2-B5-D09-VD-073,muc=VD,loigiai={$u_1=4$, $u_{n+1}=u_n+3$ nên $u_n=3n+1$; $u_{33}=100$.}]{Bạn Hải dùng các que diêm giống nhau xếp thành một dãy các hình vuông liền nhau trên mặt bàn: hình thứ nhất là $1$ hình vuông cần $4$ que; mỗi lần xếp thêm một hình vuông kề bên phải hình vừa xếp, hai hình vuông dùng chung một cạnh nên chỉ cần thêm $3$ que. Gọi $u_n$ là số que diêm cần dùng để xếp được dãy gồm $n$ hình vuông. Số que diêm cần dùng để xếp dãy gồm $33$ hình vuông là}
  {$100$.}
  {$101$.}{$132$.}{$102$.}

\cauTN[id=T11-C2-B5-D09-VD-074,muc=VD,loigiai={$u_n=4+2n$; $u_{8}=20$.}]{Đầu năm học, bạn Mai có sẵn $4$ trăm nghìn đồng trong ống heo tiết kiệm. Kể từ tháng thứ nhất, cuối mỗi tháng bạn bỏ thêm vào ống heo $2$ trăm nghìn đồng tiền tiết kiệm từ tiền ăn sáng và không rút ra lần nào. Gọi $u_n$ (trăm nghìn đồng) là số tiền có trong ống heo sau $n$ tháng. Sau $8$ tháng, số tiền trong ống heo là}
  {$20$ trăm nghìn đồng.}
  {$18$ trăm nghìn đồng.}{$16$ trăm nghìn đồng.}{$22$ trăm nghìn đồng.}

\cauTN[id=T11-C2-B5-D09-VD-075,muc=VD,loigiai={$u_1=4$, $u_{n+1}=u_n+3$ nên $u_n=3n+1$; $u_{27}=82$.}]{Bạn An dùng các que diêm giống nhau xếp thành một dãy các hình vuông liền nhau trên mặt bàn: hình thứ nhất là $1$ hình vuông cần $4$ que; mỗi lần xếp thêm một hình vuông kề bên phải hình vừa xếp, hai hình vuông dùng chung một cạnh nên chỉ cần thêm $3$ que. Gọi $u_n$ là số que diêm cần dùng để xếp được dãy gồm $n$ hình vuông. Số que diêm cần dùng để xếp dãy gồm $27$ hình vuông là}
  {$82$.}
  {$83$.}{$85$.}{$84$.}

\cauTN[id=T11-C2-B5-D09-VD-076,muc=VD,loigiai={$u_n=3+n$; $u_{11}=14$.}]{Đầu năm học, bạn Hương có sẵn $3$ trăm nghìn đồng trong ống heo tiết kiệm. Kể từ tháng thứ nhất, cuối mỗi tháng bạn bỏ thêm vào ống heo $1$ trăm nghìn đồng tiền tiết kiệm từ tiền ăn sáng và không rút ra lần nào. Gọi $u_n$ (trăm nghìn đồng) là số tiền có trong ống heo sau $n$ tháng. Sau $11$ tháng, số tiền trong ống heo là}
  {$14$ trăm nghìn đồng.}
  {$13$ trăm nghìn đồng.}{$11$ trăm nghìn đồng.}{$15$ trăm nghìn đồng.}

\cauTN[id=T11-C2-B5-D09-VD-077,muc=VD,loigiai={$u_1=4$, $u_{n+1}=u_n+3$ nên $u_n=3n+1$; $u_{37}=112$.}]{Bạn Tuấn dùng các que diêm giống nhau xếp thành một dãy các hình vuông liền nhau trên mặt bàn: hình thứ nhất là $1$ hình vuông cần $4$ que; mỗi lần xếp thêm một hình vuông kề bên phải hình vừa xếp, hai hình vuông dùng chung một cạnh nên chỉ cần thêm $3$ que. Gọi $u_n$ là số que diêm cần dùng để xếp được dãy gồm $n$ hình vuông. Số que diêm cần dùng để xếp dãy gồm $37$ hình vuông là}
  {$112$.}
  {$148$.}{$113$.}{$110$.}

\cauTN[id=T11-C2-B5-D09-VD-078,muc=VD,loigiai={$u_n=5+2n$; $u_{3}=11$.}]{Đầu năm học, bạn Hoa có sẵn $5$ trăm nghìn đồng trong ống heo tiết kiệm. Kể từ tháng thứ nhất, cuối mỗi tháng bạn bỏ thêm vào ống heo $2$ trăm nghìn đồng tiền tiết kiệm từ tiền ăn sáng và không rút ra lần nào. Gọi $u_n$ (trăm nghìn đồng) là số tiền có trong ống heo sau $n$ tháng. Sau $3$ tháng, số tiền trong ống heo là}
  {$11$ trăm nghìn đồng.}
  {$9$ trăm nghìn đồng.}{$6$ trăm nghìn đồng.}{$13$ trăm nghìn đồng.}

\cauTN[id=T11-C2-B5-D09-VD-079,muc=VD,loigiai={$u_1=4$, $u_{n+1}=u_n+3$ nên $u_n=3n+1$; $u_{15}=46$.}]{Bạn Hà dùng các que diêm giống nhau xếp thành một dãy các hình vuông liền nhau trên mặt bàn: hình thứ nhất là $1$ hình vuông cần $4$ que; mỗi lần xếp thêm một hình vuông kề bên phải hình vừa xếp, hai hình vuông dùng chung một cạnh nên chỉ cần thêm $3$ que. Gọi $u_n$ là số que diêm cần dùng để xếp được dãy gồm $n$ hình vuông. Số que diêm cần dùng để xếp dãy gồm $15$ hình vuông là}
  {$46$.}
  {$49$.}{$47$.}{$44$.}

\cauTN[id=T11-C2-B5-D09-VD-080,muc=VD,loigiai={$u_n=3+2n$; $u_{7}=17$.}]{Đầu năm học, bạn An có sẵn $3$ trăm nghìn đồng trong ống heo tiết kiệm. Kể từ tháng thứ nhất, cuối mỗi tháng bạn bỏ thêm vào ống heo $2$ trăm nghìn đồng tiền tiết kiệm từ tiền ăn sáng và không rút ra lần nào. Gọi $u_n$ (trăm nghìn đồng) là số tiền có trong ống heo sau $n$ tháng. Sau $7$ tháng, số tiền trong ống heo là}
  {$17$ trăm nghìn đồng.}
  {$15$ trăm nghìn đồng.}{$14$ trăm nghìn đồng.}{$19$ trăm nghìn đồng.}

\cauTN[id=T11-C2-B5-D09-VD-081,muc=VD,loigiai={$u_1=4$, $u_{n+1}=u_n+3$ nên $u_n=3n+1$; $u_{14}=43$.}]{Bạn Quân dùng các que diêm giống nhau xếp thành một dãy các hình vuông liền nhau trên mặt bàn: hình thứ nhất là $1$ hình vuông cần $4$ que; mỗi lần xếp thêm một hình vuông kề bên phải hình vừa xếp, hai hình vuông dùng chung một cạnh nên chỉ cần thêm $3$ que. Gọi $u_n$ là số que diêm cần dùng để xếp được dãy gồm $n$ hình vuông. Số que diêm cần dùng để xếp dãy gồm $14$ hình vuông là}
  {$43$.}
  {$44$.}{$56$.}{$46$.}

\cauTN[id=T11-C2-B5-D09-VD-082,muc=VD,loigiai={$u_n=2+\dfrac{3}{2}n$; $u_{3}=\dfrac{13}{2}$.}]{Đầu năm học, bạn Hải có sẵn $2$ trăm nghìn đồng trong ống heo tiết kiệm. Kể từ tháng thứ nhất, cuối mỗi tháng bạn bỏ thêm vào ống heo $\dfrac{3}{2}$ trăm nghìn đồng tiền tiết kiệm từ tiền ăn sáng và không rút ra lần nào. Gọi $u_n$ (trăm nghìn đồng) là số tiền có trong ống heo sau $n$ tháng. Sau $3$ tháng, số tiền trong ống heo là}
  {$\dfrac{13}{2}$ trăm nghìn đồng.}
  {$5$ trăm nghìn đồng.}{$\dfrac{9}{2}$ trăm nghìn đồng.}{$8$ trăm nghìn đồng.}

\cauDS[id=T11-C2-B5-D09-VD-083,muc=VD,loigiai={$u_n=3n+1$; $u_{33}=100$; $3n+1\le176\Leftrightarrow n\le58$. a) Đúng. b) Đúng. c) Sai. d) Đúng.}]{Bạn Đức dùng các que diêm giống nhau xếp thành một dãy các hình vuông liền nhau trên mặt bàn: hình thứ nhất là $1$ hình vuông cần $4$ que; mỗi lần xếp thêm một hình vuông kề bên phải hình vừa xếp, hai hình vuông dùng chung một cạnh nên chỉ cần thêm $3$ que. Gọi $u_n$ là số que diêm cần dùng để xếp được dãy gồm $n$ hình vuông.}{%
  \y{$u_2=7$.}{d}
  \y{$u_{n+1}=u_n+3$ với mọi $n\ge1$.}{d}
  \y{$u_{33}=103$.}{s}
  \y{Với $176$ que diêm, bạn ấy xếp được tối đa $58$ hình vuông.}{d}}

\cauDS[id=T11-C2-B5-D09-VD-084,muc=VD,loigiai={$u_n=3n+1$; $u_{5}=16$; $3n+1\le139\Leftrightarrow n\le46$. a) Đúng. b) Đúng. c) Sai. d) Sai.}]{Bạn Chi dùng các que diêm giống nhau xếp thành một dãy các hình vuông liền nhau trên mặt bàn: hình thứ nhất là $1$ hình vuông cần $4$ que; mỗi lần xếp thêm một hình vuông kề bên phải hình vừa xếp, hai hình vuông dùng chung một cạnh nên chỉ cần thêm $3$ que. Gọi $u_n$ là số que diêm cần dùng để xếp được dãy gồm $n$ hình vuông.}{%
  \y{$u_2=7$.}{d}
  \y{$u_{n+1}=u_n+3$ với mọi $n\ge1$.}{d}
  \y{$u_{5}=19$.}{s}
  \y{Với $139$ que diêm, bạn ấy xếp được tối đa $47$ hình vuông.}{s}}

\cauDS[id=T11-C2-B5-D09-VD-085,muc=VD,loigiai={$u_n=3n+1$; $u_{13}=40$; $3n+1\le72\Leftrightarrow n\le23$. a) Đúng. b) Đúng. c) Đúng. d) Sai.}]{Bạn Hà dùng các que diêm giống nhau xếp thành một dãy các hình vuông liền nhau trên mặt bàn: hình thứ nhất là $1$ hình vuông cần $4$ que; mỗi lần xếp thêm một hình vuông kề bên phải hình vừa xếp, hai hình vuông dùng chung một cạnh nên chỉ cần thêm $3$ que. Gọi $u_n$ là số que diêm cần dùng để xếp được dãy gồm $n$ hình vuông.}{%
  \y{$u_2=7$.}{d}
  \y{$u_{n+1}=u_n+3$ với mọi $n\ge1$.}{d}
  \y{$u_{13}=40$.}{d}
  \y{Với $72$ que diêm, bạn ấy xếp được tối đa $24$ hình vuông.}{s}}

\cauDS[id=T11-C2-B5-D09-VD-086,muc=VD,loigiai={$u_n=3n+1$; $u_{36}=109$; $3n+1\le137\Leftrightarrow n\le45$. a) Đúng. b) Đúng. c) Đúng. d) Đúng.}]{Bạn Nam dùng các que diêm giống nhau xếp thành một dãy các hình vuông liền nhau trên mặt bàn: hình thứ nhất là $1$ hình vuông cần $4$ que; mỗi lần xếp thêm một hình vuông kề bên phải hình vừa xếp, hai hình vuông dùng chung một cạnh nên chỉ cần thêm $3$ que. Gọi $u_n$ là số que diêm cần dùng để xếp được dãy gồm $n$ hình vuông.}{%
  \y{$u_2=7$.}{d}
  \y{$u_{n+1}=u_n+3$ với mọi $n\ge1$.}{d}
  \y{$u_{36}=109$.}{d}
  \y{Với $137$ que diêm, bạn ấy xếp được tối đa $45$ hình vuông.}{d}}

\cauTLN[id=T11-C2-B5-D09-VD-087,muc=VD,loigiai={$u_n=3n+1\le186\Leftrightarrow n\le61$.}]{Bạn Hương dùng các que diêm giống nhau xếp thành một dãy các hình vuông liền nhau trên mặt bàn: hình thứ nhất là $1$ hình vuông cần $4$ que; mỗi lần xếp thêm một hình vuông kề bên phải hình vừa xếp, hai hình vuông dùng chung một cạnh nên chỉ cần thêm $3$ que. Gọi $u_n$ là số que diêm cần dùng để xếp được dãy gồm $n$ hình vuông. Với $186$ que diêm, bạn ấy xếp được dãy có nhiều nhất bao nhiêu hình vuông?}{61}

\cauTLN[id=T11-C2-B5-D09-VD-088,muc=VD,loigiai={$5+\dfrac{1}{2}n\ge10\Leftrightarrow n\ge10$; chọn $n=10$.}]{Đầu năm học, bạn Phúc có sẵn $5$ trăm nghìn đồng trong ống heo tiết kiệm. Kể từ tháng thứ nhất, cuối mỗi tháng bạn bỏ thêm vào ống heo $\dfrac{1}{2}$ trăm nghìn đồng tiền tiết kiệm từ tiền ăn sáng và không rút ra lần nào. Gọi $u_n$ (trăm nghìn đồng) là số tiền có trong ống heo sau $n$ tháng. Sau ít nhất bao nhiêu tháng thì số tiền trong ống heo đạt từ $10$ trăm nghìn đồng trở lên?}{10}

\cauTLN[id=T11-C2-B5-D09-VD-089,muc=VD,loigiai={$u_n=3n+1\le241\Leftrightarrow n\le80$.}]{Bạn Khoa dùng các que diêm giống nhau xếp thành một dãy các hình vuông liền nhau trên mặt bàn: hình thứ nhất là $1$ hình vuông cần $4$ que; mỗi lần xếp thêm một hình vuông kề bên phải hình vừa xếp, hai hình vuông dùng chung một cạnh nên chỉ cần thêm $3$ que. Gọi $u_n$ là số que diêm cần dùng để xếp được dãy gồm $n$ hình vuông. Với $241$ que diêm, bạn ấy xếp được dãy có nhiều nhất bao nhiêu hình vuông?}{80}

\cauTLN[id=T11-C2-B5-D09-VD-090,muc=VD,loigiai={$3+\dfrac{3}{2}n\ge38\Leftrightarrow n\ge\dfrac{70}{3}$; chọn $n=24$.}]{Đầu năm học, bạn Đức có sẵn $3$ trăm nghìn đồng trong ống heo tiết kiệm. Kể từ tháng thứ nhất, cuối mỗi tháng bạn bỏ thêm vào ống heo $\dfrac{3}{2}$ trăm nghìn đồng tiền tiết kiệm từ tiền ăn sáng và không rút ra lần nào. Gọi $u_n$ (trăm nghìn đồng) là số tiền có trong ống heo sau $n$ tháng. Sau ít nhất bao nhiêu tháng thì số tiền trong ống heo đạt từ $38$ trăm nghìn đồng trở lên?}{24}

\cauTN[id=T11-C2-B5-D09-VDC-091,muc=VDC,loigiai={$u_1=4$, $u_{n+1}=u_n+3$ nên $u_n=3n+1$; $u_{10}=31$.}]{Bạn Cường dùng các que diêm giống nhau xếp thành một dãy các hình vuông liền nhau trên mặt bàn: hình thứ nhất là $1$ hình vuông cần $4$ que; mỗi lần xếp thêm một hình vuông kề bên phải hình vừa xếp, hai hình vuông dùng chung một cạnh nên chỉ cần thêm $3$ que. Gọi $u_n$ là số que diêm cần dùng để xếp được dãy gồm $n$ hình vuông. Số que diêm cần dùng để xếp dãy gồm $10$ hình vuông là}
  {$31$.}
  {$29$.}{$30$.}{$32$.}

\cauTN[id=T11-C2-B5-D09-VDC-092,muc=VDC,loigiai={$u_n=3+\dfrac{1}{2}n$; $u_{11}=\dfrac{17}{2}$.}]{Đầu năm học, bạn Sơn có sẵn $3$ trăm nghìn đồng trong ống heo tiết kiệm. Kể từ tháng thứ nhất, cuối mỗi tháng bạn bỏ thêm vào ống heo $\dfrac{1}{2}$ trăm nghìn đồng tiền tiết kiệm từ tiền ăn sáng và không rút ra lần nào. Gọi $u_n$ (trăm nghìn đồng) là số tiền có trong ống heo sau $n$ tháng. Sau $11$ tháng, số tiền trong ống heo là}
  {$\dfrac{17}{2}$ trăm nghìn đồng.}
  {$8$ trăm nghìn đồng.}{$\dfrac{11}{2}$ trăm nghìn đồng.}{$9$ trăm nghìn đồng.}

\cauTN[id=T11-C2-B5-D09-VDC-093,muc=VDC,loigiai={Tính lần lượt các số hạng: $u_{2}=100$ con.}]{Một trang trại nuôi cá có $100$ con cá giống ở đầu tháng thứ nhất. Sau mỗi tháng, số cá trong hồ tăng thêm $20\%$ so với đầu tháng do sinh sản, đồng thời cuối mỗi tháng trang trại bán ra đúng $20$ con. Gọi $u_n$ là số cá trong hồ sau $n$ tháng ($u_0=100$), ta có $u_{n+1}=\dfrac{6}{5} u_n-20$. Số cá trong hồ sau $2$ tháng là}
  {$100$.}
  {$101$.}{$120$.}{$99$.}

\cauTN[id=T11-C2-B5-D09-VDC-094,muc=VDC,loigiai={$u_1=4$, $u_{n+1}=u_n+3$ nên $u_n=3n+1$; $u_{33}=100$.}]{Bạn Linh dùng các que diêm giống nhau xếp thành một dãy các hình vuông liền nhau trên mặt bàn: hình thứ nhất là $1$ hình vuông cần $4$ que; mỗi lần xếp thêm một hình vuông kề bên phải hình vừa xếp, hai hình vuông dùng chung một cạnh nên chỉ cần thêm $3$ que. Gọi $u_n$ là số que diêm cần dùng để xếp được dãy gồm $n$ hình vuông. Số que diêm cần dùng để xếp dãy gồm $33$ hình vuông là}
  {$100$.}
  {$102$.}{$103$.}{$101$.}

\cauTN[id=T11-C2-B5-D09-VDC-095,muc=VDC,loigiai={$u_n=3+\dfrac{3}{2}n$; $u_{12}=21$.}]{Đầu năm học, bạn Mai có sẵn $3$ trăm nghìn đồng trong ống heo tiết kiệm. Kể từ tháng thứ nhất, cuối mỗi tháng bạn bỏ thêm vào ống heo $\dfrac{3}{2}$ trăm nghìn đồng tiền tiết kiệm từ tiền ăn sáng và không rút ra lần nào. Gọi $u_n$ (trăm nghìn đồng) là số tiền có trong ống heo sau $n$ tháng. Sau $12$ tháng, số tiền trong ống heo là}
  {$21$ trăm nghìn đồng.}
  {$\dfrac{39}{2}$ trăm nghìn đồng.}{$18$ trăm nghìn đồng.}{$\dfrac{45}{2}$ trăm nghìn đồng.}

\cauTN[id=T11-C2-B5-D09-VDC-096,muc=VDC,loigiai={Tính lần lượt các số hạng: $u_{2}=221$ con.}]{Một trang trại nuôi cá có $200$ con cá giống ở đầu tháng thứ nhất. Sau mỗi tháng, số cá trong hồ tăng thêm $10\%$ so với đầu tháng do sinh sản, đồng thời cuối mỗi tháng trang trại bán ra đúng $10$ con. Gọi $u_n$ là số cá trong hồ sau $n$ tháng ($u_0=200$), ta có $u_{n+1}=\dfrac{11}{10} u_n-10$. Số cá trong hồ sau $2$ tháng là}
  {$221$.}
  {$222$.}{$231$.}{$242$.}

\cauDS[id=T11-C2-B5-D09-VDC-097,muc=VDC,loigiai={$u_n=3n+1$; $u_{21}=64$; $3n+1\le179\Leftrightarrow n\le59$. a) Đúng. b) Đúng. c) Đúng. d) Đúng.}]{Bạn Đức dùng các que diêm giống nhau xếp thành một dãy các hình vuông liền nhau trên mặt bàn: hình thứ nhất là $1$ hình vuông cần $4$ que; mỗi lần xếp thêm một hình vuông kề bên phải hình vừa xếp, hai hình vuông dùng chung một cạnh nên chỉ cần thêm $3$ que. Gọi $u_n$ là số que diêm cần dùng để xếp được dãy gồm $n$ hình vuông.}{%
  \y{$u_2=7$.}{d}
  \y{$u_{n+1}=u_n+3$ với mọi $n\ge1$.}{d}
  \y{$u_{21}=64$.}{d}
  \y{Với $179$ que diêm, bạn ấy xếp được tối đa $59$ hình vuông.}{d}}

\cauDS[id=T11-C2-B5-D09-VDC-098,muc=VDC,loigiai={$u_n=3n+1$; $u_{10}=31$; $3n+1\le139\Leftrightarrow n\le46$. a) Đúng. b) Đúng. c) Đúng. d) Sai.}]{Bạn Nam dùng các que diêm giống nhau xếp thành một dãy các hình vuông liền nhau trên mặt bàn: hình thứ nhất là $1$ hình vuông cần $4$ que; mỗi lần xếp thêm một hình vuông kề bên phải hình vừa xếp, hai hình vuông dùng chung một cạnh nên chỉ cần thêm $3$ que. Gọi $u_n$ là số que diêm cần dùng để xếp được dãy gồm $n$ hình vuông.}{%
  \y{$u_2=7$.}{d}
  \y{$u_{n+1}=u_n+3$ với mọi $n\ge1$.}{d}
  \y{$u_{10}=31$.}{d}
  \y{Với $139$ que diêm, bạn ấy xếp được tối đa $47$ hình vuông.}{s}}

\cauTLN[id=T11-C2-B5-D09-VDC-099,muc=VDC,loigiai={$u_n=3n+1\le393\Leftrightarrow n\le130$.}]{Bạn Quỳnh dùng các que diêm giống nhau xếp thành một dãy các hình vuông liền nhau trên mặt bàn: hình thứ nhất là $1$ hình vuông cần $4$ que; mỗi lần xếp thêm một hình vuông kề bên phải hình vừa xếp, hai hình vuông dùng chung một cạnh nên chỉ cần thêm $3$ que. Gọi $u_n$ là số que diêm cần dùng để xếp được dãy gồm $n$ hình vuông. Với $393$ que diêm, bạn ấy xếp được dãy có nhiều nhất bao nhiêu hình vuông?}{130}

\cauTLN[id=T11-C2-B5-D09-VDC-100,muc=VDC,loigiai={$5+\dfrac{3}{2}n\ge27\Leftrightarrow n\ge\dfrac{44}{3}$; chọn $n=15$.}]{Đầu năm học, bạn Mai có sẵn $5$ trăm nghìn đồng trong ống heo tiết kiệm. Kể từ tháng thứ nhất, cuối mỗi tháng bạn bỏ thêm vào ống heo $\dfrac{3}{2}$ trăm nghìn đồng tiền tiết kiệm từ tiền ăn sáng và không rút ra lần nào. Gọi $u_n$ (trăm nghìn đồng) là số tiền có trong ống heo sau $n$ tháng. Sau ít nhất bao nhiêu tháng thì số tiền trong ống heo đạt từ $27$ trăm nghìn đồng trở lên?}{15}
